% ===== §1 The Plain Happiness across Theoretical Horizons =====
\section{The Plain Happiness across Theoretical Horizons}
\label{p19:sec:theories}

\noindent
The phenomenon having been delineated, this chapter asks how the major theoretical traditions explain it. The method is comparative and expository: each tradition is taken on its own terms and allowed to illuminate the plain happiness once, in the idiom proper to it, without being corrected or absorbed into any other. The traditions are placed side by side rather than woven together, and the work of drawing them into relation---of asking what they share and where they fail in common---is reserved for the final section, \xref{p19:sub:theories-convergence}. The aim of the juxtaposition is not eclecticism but exposure: the thickness of the phenomenon is best shown by how many independent traditions find in it something they recognise as their own. Throughout, points requiring exact textual attribution are flagged for verification rather than asserted from memory.

A word on order. The sections proceed, loosely, from the most proximate register---the bare structure of two people attending to the same thing---outward to the most structural, the place of such happiness in an economy of production. The ordering is heuristic; the sections are mutually independent and may be read in any sequence.

\subsection{Joint attention and intersubjectivity}
\label{p19:sub:theories-joint}

The minimal structure of the datum, considered psychologically, is that two subjects attend together to the same object without the need for speech. The question this tradition poses is why such a wordless sharing of orientation should itself be a form of contentment, rather than a mere coordination of two perceptual states.

Developmental psychology has made the sharing of attention foundational to its account of how human sociality arises. In the framework of \textcite{trevarthen1979}, a \emph{primary} intersubjectivity---the early affective attunement of infant and caregiver, a direct coordination of two subjectivities---is succeeded, after roughly the ninth month, by a \emph{secondary} intersubjectivity in which a third term, the shared object, is drawn into the dyad, so that the relation acquires a triadic structure: I, you, and the object attended to together. The plain happiness of watching the rain together has the form of a mature secondary intersubjectivity: a triadic engagement in which the object is not a topic of communication but the common pole of a shared regard.

\textcite{tomasello2008} makes the sharing of intentionality the cognitive origin of human cooperation and culture, and isolates as its core a capacity for \emph{joint} attention that is not reducible to two parallel acts of individual attention: there is a ``we'' that is not the sum of two ``I''s, a shared intentional state with its own structure. On this view the togetherness of the datum is not an aggregate but a unit; the happiness, if it is located in the togetherness, is located in something genuinely irreducible to the two parties taken separately.

\textcite{stern1985} adds a feature especially pertinent here. In his account of \emph{affect attunement}, what passes between infant and caregiver in their wordless exchanges is not the content of a feeling but its dynamic contour---its \emph{vitality affect}, the shape of its rising and falling. This is a description of a sharing that has no objective content and yet a definite dynamic form: precisely the combination---contentless yet moving---that the datum exhibits in its ``faint movement, not inertness.'' The tradition thus already possesses a vocabulary for a togetherness that is full of movement without being full of any particular content.

\paragraph{What it illuminates.} The minimal unit of the plain happiness of being-with is a wordless \emph{co-orientation}, a shared direction of regard, rather than any communicated content. Watching the rain together is satisfying because it instantiates the purest case of a shared attentional state: a being-present-together that requires nothing to be said.

\paragraph{Its limit.} In the developmental framework, joint attention is positioned as a cognitive \emph{milestone}, or a capacity, rather than as a mode of happiness in its own right: it explains how the ``we'' becomes possible, not why the ``we'' should be a fullness. More tellingly, the teleological cast of the framework tends to read joint attention as a way-station \emph{toward} language and cooperation---a means---whereas the plain happiness of watching the rain together is precisely an end in itself, leading on to no subsequent attainment. This tension---between the milestone read as means and the togetherness lived as end in itself---is one the tradition cannot dissolve from within.

\subsection{Psychoanalysis: drive, the symbolic--real gap, and being quietly present}
\label{p19:sub:theories-psych}

Psychoanalysis takes drive and desire as its underlying dynamics, and this gives the datum a double edge. On the one hand, if happiness is conceived as the discharge or satisfaction of desire, then an uneventful, low-tension being-together is hard to count as a fullness rather than as the mere temporary absence of tension. On the other hand---and this is the feature the present paper will draw upon most heavily later---psychoanalysis is almost alone among the traditions surveyed here in offering a candidate \emph{source} for a perpetually renewed dynamism: if such a source can be identified, the tradition will have supplied not merely a picture of the plain happiness but an engine beneath it. The first edge is developed in this chapter; the constitutive use of the second is deferred to the next.

\paragraph{Freud: drive and the economy of desire.} Freud's \emph{Trieb} is a frontier concept between the somatic and the psychical, possessing a source, a pressure, an aim, and an object, of which the object is the most variable while the pressure is constant \parencite{freud1915}. The drive is thus, by its structure, never finally satisfied: the object may be replaced, the pressure does not extinguish. This is the earliest psychoanalytic model of a dynamism that does not stop. Against it stands a contrary tendency in the same corpus: the \emph{Nirvana principle}, the trend toward the reduction of excitation to zero, which Freud (adopting Barbara Low's term) formulated in \emph{Beyond the Pleasure Principle} and aligned with the death drive---distinct from, though entangled with, the \emph{constancy principle}'s tendency to a low or stable level \parencite{freud1920}. Under it the plain happiness would be read as a tending-toward-stillness close to inertness---the psychoanalytic form of precisely the ``empty stillness'' the present paper must guard against. The fullness of the low-tension condition has therefore to be distinguished, within psychoanalysis itself, from its deathlike counterfeit.

\paragraph{Lacan: desire as the gap between the symbolic and the real.} The decisive resource lies here. In Lacan's topology the real is what resists symbolisation, and the symbolic, the order of signifiers, never captures it without remainder \parencite{lacan2006}. This structural shortfall of the symbolic with respect to the real---the irreducible gap between them---is what sets desire in motion; the \emph{objet petit a} marks the leftover of that gap, the object-cause around which desire turns. Desire is therefore never satisfied, perpetually displaced along the chain of signifiers, never filled by any particular object, not because it lacks some specific thing but because it is driven by the unclosable gap itself.

The proposition to be carried forward is this: \emph{the gap between the symbolic and the real is a source-model for a generative dynamism}. It accounts for why there should be an unceasing movement, rather than positing that movement as an unexplained axiom. Generation does not stop, on this account, because the gap cannot be closed.

But a tension is built into the resource, and it must be marked here rather than resolved. In Lacan the gap drives the \emph{perpetual frustration} of desire---an economy of lack, of the \emph{manque-à-être}---in which the subject, driven by the gap, pursues and forever misses. The datum of this paper, by contrast, is a \emph{plain, contented, low-tension} happiness. How can a dynamism driven by an unclosable gap, condemned to miss its object, be the source of a contented fullness? The contradiction is left standing here, suspended as the entry-point of the next chapter.

\paragraph{Winnicott: the capacity to be alone---the object-relations wing.} A different wing of psychoanalysis describes the plain togetherness directly. Winnicott's paradoxical formulation is the \emph{capacity to be alone in the presence of another} \parencite{winnicott1958}: the mature ability to be alone is founded on an early experience of being quietly alone in the mother's presence, where no interaction and no tension of desire is required, while an unobtrusive ego-relatedness is sustained. This is almost an archetype of the plain happiness of being-with: the adult form of being, quietly, oneself in the other's presence. The associated notion of \emph{going-on-being}, whose opposite is impingement, names just such an uninterrupted, plain continuance of existence. Yet this wing stands outside the economy of drive, and so cannot answer the question of the dynamism's source; while the wing that can supply the source (Lacan) locks it in lack. The two wings hold opposite ends, and a later chapter must rejoin them: the dynamism of the drive and the quiet of object-relatedness, reunited in a single account.

\paragraph{Bion and extensions.} The picture is filled out by Bion's \emph{reverie} and the container--contained relation, in which wordless accompaniment is itself a structure of containing \parencite{bion1962}, and by Balint's account of a harmonious, conflict-free interpenetration, a benign regression in which the plain being-together approaches a quiet commingling \parencite{balint1968}.

\paragraph{What it illuminates.} The plain happiness is a being-together that needs no tension of desire to sustain it (Winnicott); the highest intimacy may be one in which ego-relatedness is quietly present without the activation of drive. And, uniquely among the traditions here, psychoanalysis touches what lies \emph{beneath} the plain happiness---the question of where its dynamism comes from---and offers a candidate answer in the symbolic--real gap.

\paragraph{Its limit.} The dominant lines (the Freudian Nirvana principle, Lacanian desire) tend to read the plain happiness as the temporary absence of tension, a subtraction, or to lock its dynamism into an economy of lack. The Winnicottian line comes closest to a positive account but, standing outside the economy of drive, does not address the question of the source; its ``uninterrupted'' is a negative determination, not yet a positive generation. Psychoanalysis thus bequeaths to the sequel two legacies: an \emph{asset}---the symbolic--real gap as engine of dynamism---and a \emph{debt}---the default confinement of that dynamism to the mode of lack. The task of the next chapter is to inherit the asset while breaking with the debt.

\subsection{Phenomenology: passive synthesis, intercorporeity, releasement}
\label{p19:sub:theories-phenom}

Phenomenology's commitment is to describe how things are given. The question it poses to the datum is whether the still afternoon, as a given fullness, can be grasped \emph{at the level of description}---whether one can characterise a ``full without event'' in terms of the structure of its givenness.

\paragraph{Husserl: passive synthesis and inner time.} For Husserl the fullness of the afternoon need not be constituted by any active synthesis; at the level of \emph{passive synthesis} it is already given as replete \parencite{husserl2001}. And the structure of inner time-consciousness---the threefold of retention, primal impression, and protention---means that the present is never a static point but always carries, in its protentional edge, a slight openness toward what is coming. The still afternoon is therefore not a motionless instant but a continually self-renewing flow: phenomenology, at the descriptive level, already registers a faint movement within the stillness. The fifth Cartesian Meditation's account of \emph{pairing} (\emph{Paarung}), in which intersubjectivity is passively constituted in the coupling of bodies, supplies in turn a descriptive ground for the joint attention of \xref{p19:sub:theories-joint}.

\paragraph{Merleau-Ponty: intercorporeity and the flesh of the shared.} For Merleau-Ponty the watching-together is two bodies sharing one world; the being-with is given pre-reflectively in \emph{intercorporeity}, the intertwining (\emph{chiasme}) of the flesh, rather than in any agreement of two consciousnesses \parencite{merleau-ponty1962}. The ``we'' of the plain happiness is a matter of the flesh, not of a contract of minds.

\paragraph{Heidegger: attunement and releasement.} The plain happiness is a \emph{Stimmung}, an attunement prior to the division of subject and object, a low-intensity, open mode of finding-oneself-disposed; it is not the absence of mood but a quiet, open one. \emph{Gelassenheit}---releasement, the letting-be that neither forces nor manipulates but lets things be as they are \parencite{heidegger1966}---names the ontological vindication of the afternoon's ``doing nothing,'' and stands in close kinship with the Daoist non-doing. \emph{Mitsein} ensures that being-with is an original structure, not a later addition.

\paragraph{What it illuminates.} The plain happiness is a \emph{given fullness}, with an attunement-structure (\emph{Stimmung}) and a faintly open temporal structure (protention); phenomenology lets us say that it is not empty but a replete givenness.

\paragraph{Its limit.} Phenomenology is expert at describing how a thing is given and \emph{stops at description}. It can render the structure of the afternoon's givenness with great precision, yet methodologically brackets the genetic-dynamic question of why this is a \emph{generation}, whence this fullness continually grows. Protention touches the ``faint movement within stillness,'' but phenomenology does not ask after the \emph{source} of that movement, nor after its plasticity under historical and material conditions.

\subsection{The eudaimonist tradition}
\label{p19:sub:theories-eudaimonia}

The question here is in what sense the classical theories of well-being can accommodate a happiness that is plain yet full. The decisive fault-line runs between conceiving happiness as an \emph{activity} and conceiving it as a \emph{state}, and between conceiving pleasure as \emph{kinetic} (in process) and as \emph{static} (a settled absence of disturbance). The plain happiness falls exactly on this line: uneventful (near the static), yet full (near the activity).

\paragraph{Aristotle: \emph{eudaimonia} as \emph{energeia}.} For Aristotle happiness is not a state but the soul's activity in accordance with virtue---an \emph{energeia} \parencite{aristotle2000}. Its highest form, \emph{theoria}, is an activity outwardly still yet inwardly most active, which gives the ``plain yet full'' a classical foothold; and \emph{scholē}, leisure, is not idleness but the condition under which free activity becomes possible. Watching the rain together might be read as a low-intensity shared \emph{energeia} rather than as doing nothing. The strain is that Aristotelian \emph{theoria} is solitary and self-sufficient, ill-suited to the dimension of being-with, and that it retains a determinate object, whereas the plain being-together has none.

\paragraph{Epicurus: katastematic pleasure.} The highest pleasure is \emph{katastematic}---not the kinetic pleasure of sensory stimulation but the settled state of freedom from pain and disturbance (\emph{aponia}, \emph{ataraxia}) \parencite{epicurus1994}. This is nearly the most direct classical advocate of ``plainness is fullness'': the plain is not the absence of pleasure but the very form of the highest pleasure. Yet the katastematic is a subtractive ideal---the removal of want and disturbance---whose fullness is a ``being-without-deficit'' rather than a ``being-that-grows''; it cannot speak to the growing side of the plain happiness.

\paragraph{Stoicism: \emph{apatheia} and accord with nature.} The plain as a stability undisturbed by the passions, the product of accord with cosmic reason---an achievement, not numbness. But \emph{apatheia} remains a negative determination, a freedom-from-disturbance, and its ideal of individual self-sufficiency is likewise ill-fitted to the growth of a shared condition.

\paragraph{Eastern parallels: Zhuangzi's plainness and the Confucian joy.} Zhuangzi's \emph{tiandan} (plainness), the ``empty room that gives forth light'' (\emph{xushi sheng bai}), and the saying that the plain and limitless is that from which every beauty follows---here plainness is the \emph{condition of generation} (emptied, it lets something be born), which is the crucial difference from the Western static-pleasure tradition: Zhuangzi's plainness is generative rather than merely undisturbed \parencite{zhuangzi2009}. The Confucian joy of Yan Hui amid poverty, and the ``Zeng Dian'' vision of bathing in the river and returning singing, place a plain, shared, seasonal joy at the heart of virtue \parencite{confucius1998}. The Eastern tradition, Zhuangzi above all, already harbours the ``stillness that generates'' the Western eudaimonism lacks---an important anticipation of the next chapter, though it is left there as an anticipation, untheorised.

\paragraph{What it illuminates and its limit.} The eudaimonist tradition wins for the plain a positive evaluative standing (it is not a deficiency). But the Western mainstream (Epicurus, the Stoics) reads it as a subtractive, undisturbed stillness; Aristotle's \emph{energeia} has a ``motion within stillness'' but confines it to solitary contemplation with a determinate object, ill-suited to the growth of an objectless being-with. Only the Eastern line (Zhuangzi) escapes this, harbouring the ``stillness that generates''---but as an intimation, not yet a theory.

\subsection{Neuroscience: low arousal and the calm of prediction}
\label{p19:sub:theories-neuro}

The neural picture of the plain happiness is, at first glance, one of ``lowness'': low arousal, low prediction error, low phasic dopaminergic activity. The question is whether neuroscience can distinguish, within this lowness, two conditions---an empty lowness and a full one. The question is only raised here; its positive resolution is deferred to \xref{p19:sec:neuro}.

\paragraph{Arousal and the autonomic tone of calm.} The plain being-together corresponds to a safe, low-arousal, parasympathetically dominant tone---in the terms of the polyvagal account, a ventral-vagal state of safety and social engagement \parencite{porges2007}, physiologically the signature of ``safe presence.''

\paragraph{Predictive processing and free energy.} The plain corresponds to a persistently low prediction error: the world is well modelled and the cortex need issue no large signal of surprise \parencite{friston2010}. But in the dominant narrative low prediction error can equally point to \emph{low information}, to boredom---and the difference between empty and full cannot be told apart at this level. This is flagged here as exactly the criterial problem that \xref{p19:sub:neuro-fullness} must address: not the height of the error but the tone of the neuromodulation and the character of the resting activity beneath it.

\paragraph{The default mode network and rest.} Resting-state activity in the default mode network is associated both with calm introspection and with rumination and negative self-reference \parencite{raichle2001}; ``rest'' is therefore in itself neutral, its value depending on the character of the activity---again pointing to the criterion of \xref{p19:sec:neuro}.

\paragraph{What it illuminates and its limit.} The plain happiness has a neural substrate, a describable condition of calm (safe, low-error, parasympathetic). But the mainstream framework tends to read ``low arousal / low error'' as \emph{low value} (boredom, low information), and cannot at its own level distinguish empty from full. More fundamentally, neuroscience describes the \emph{substrate}, and does not answer the normative--ontological question of what this fullness is or why it is an end in itself; it is a necessary witness, not a ground. The gulf between the mechanistic ``low'' and the experiential ``full'' is one neuroscience cannot itself cross.

\subsection{Structuralism and post-structuralism}
\label{p19:sub:theories-structuralism}

If meaning is fixed by position within a system of differences---with no positive terms, only differences---then what place does the plain, as a low-intensity condition of meaning, occupy semiotically? Can structuralism speak of \emph{fullness}, or only of \emph{operation}?

\paragraph{The structuralist ground.} On the Saussurean basis---\emph{langue} and \emph{parole}, value as differential, a system of differences without positive terms \parencite{saussure1959}---the plain figures as a temporary anchoring of the chain of signifiers, a low-intensity state of the play of difference.

\paragraph{Barthes: the Neutral, tenderness, the lover's discourse.} In \emph{A Lover's Discourse} the undramatic, eventless moments of tenderness; in the lecture course on \emph{the Neutral}, the celebration of a Neutral that suspends the conflict of paradigms \parencite{barthes1978,barthes2005}---the plain as a \emph{le Neutre}: a suspended presence that does not stand within the intensities of binary opposition. Here structuralism's heir comes closest to naming the plain positively, as a low-intensity mode of meaning.

\paragraph{The post-structuralist extension.} Deleuze admits a positive extension---the plain as a low-intensity becoming, a micro-perception on a plane of immanence \parencite{deleuze1987}; but Deleuze has already departed the structuralist ``difference as absence'' toward an affirmative becoming, and so constitutes a potential bridge to the next chapter. His ``becoming,'' however, is impersonal and cosmological, and lacks the relational, two-subject, self-continuing dimension this paper requires; the difference is to be marked in \xref{p19:sub:gen-process}, not here.

\paragraph{What it illuminates and its limit.} The plain is a low-intensity mode of the operation of meaning, with a definite semiotic place (the Neutral, tenderness, the anchoring of the signifier). But classical structuralism is adept at saying how meaning \emph{operates} and inept at saying how meaning is \emph{full}: on a base of ``difference, no positive terms,'' a positive fullness is nearly unsayable---it can locate the plain's ``place'' but not give its ``repletion.'' Post-structuralism (Deleuze) turns toward affirmative becoming, but its becoming is impersonal and cosmological, lacking the relational and two-subject character at issue here---a difference deferred to the next chapter.

\subsection{Political economy and Marxism}
\label{p19:sub:theories-politecon}

This section moves the plain happiness into the historical-material horizon: in a social form governed by the imperative of production and exchange-value, what place does an afternoon of doing nothing occupy? Is it waste, to be optimised away by the logic of capital, or an escape from and resistance to that logic?

\paragraph{Labour, alienation, and free time.} Against alienated labour---the fourfold estrangement of the 1844 Manuscripts---Marx sets \emph{disposable time}, free time, and ``the development of human powers as an end in itself'' \parencite{marxgrundrisse1973,marx1976}. The plain, non-productive time of being-together is the obverse of alienated labour: in it, the exercise of human powers becomes an end in itself rather than being pressed into the service of exchange-value. Marx's ``realm of freedom,'' which begins where labour determined by necessity ends, gives the plain happiness a positive location in political economy.

\paragraph{The critique of productivism and leisure.} Lafargue's polemic against the cult of labour \parencite{lafargue1883}; Adorno's analysis of the colonisation of leisure by the culture industry, in which ``free time'' is still shaped by the logic of production \parencite{adorno1991}. A tension follows: even when one escapes productive time, leisure may be recolonised by the logic of consumption---so that whether the plain happiness's ``non-productivity'' can truly escape is itself a political question.

\paragraph{Reproductive labour and invisibility.} The material conditions that make ``an afternoon of doing nothing together'' possible---who sustains the reproduction of life, by what division of labour this leisure is underwritten---cannot be romantically effaced \parencite{federici2004}. The plain happiness has material preconditions, and honesty requires acknowledging them; this connects to the questions of power and asymmetry treated in Paper~XVII \parencite{paperxvii}.

\paragraph{What it illuminates and its limit.} Political economy can locate the plain happiness's \emph{political value} (as the obverse of alienation, an escape from the imperative of production) and its \emph{material conditions}. But its determination is primarily negative and relational: the plain as the obverse of alienation, as resistance to productivism---and resistance is a negative determination, not yet a positive account of the fullness itself. It tells us what the plain is \emph{not} (not production, not alienation) and what it \emph{opposes} (productivism), but not what its repletion itself \emph{is}, nor how it continues to generate.

\subsection{Aesthetics}
\label{p19:sub:theories-aesthetics}

\paragraph{Kant.} Purposiveness without purpose (\emph{Zweckmäßigkeit ohne Zweck}), the disinterested satisfaction of the aesthetic, the free play of the cognitive faculties \parencite{kant1987}---the plain as a purposeless free play of the faculties, satisfying without serving any end. This gives the afternoon's ``doing nothing'' an aesthetic vindication: a contentment that is precisely not instrumental.

\paragraph{East-Asian aesthetics.} The Japanese \emph{ma} (the interval, the negative space) and \emph{wabi-sabi}; the Chinese painterly category of \emph{liubai} (leaving-blank) and of \emph{pingdan tianzhen} (plain and natural), in which the plain is the \emph{highest}, not the lowest, aesthetic category---the \emph{yipin} or ``untrammelled grade'' of Ni Zan and the literati tradition. That plainness ranks as the supreme aesthetic value in East Asia directly supports this paper's ``the plain is fulfilment.''

\paragraph{Its limit.} Aesthetics accords the plain its \emph{highest} evaluative standing, yet relies on a particular framework of transcendence (aesthetic autonomy), and offers no account internal to \emph{relation and history} of the plain as a generation.

\subsection{Process philosophy and relational ontology}
\label{p19:sub:theories-process}

The traditions juxtaposed in this section stand closest to the account this paper will develop---so close that, left undemarcated, they would absorb it. Their function here therefore differs from that of the others: where most were marked as ``grasping the stillness and missing the growth,'' these have already placed \emph{generation}, \emph{process}, or \emph{relation} at the centre of their ontology, and are the strongest of the ``marginal symptoms'' that escape the common blind spot of \xref{p19:sub:theories-convergence}. Yet---as the close of this section notes---they either lack the dimension of justice or take \emph{completion} as their rhythm, and so do not reach the generative relation this paper requires. The full demarcation is reserved for \xref{p19:sub:gen-process}; here they are placed alongside the rest, illuminating once.

\paragraph{Bergson: duration and the vital impulse.} The plain afternoon as pure \emph{durée}---an interpenetrating flow of experience not parcelled into a sequence of events, whose ``uneventfulness'' is precisely the qualitative repletion that resists spatialised measurement \parencite{bergson1910}. \emph{Durée} supplies a philosophy-of-time description for ``full without event.'' Its limit (deferred to \xref{p19:sub:gen-process}): \emph{durée} is the duration of an individual consciousness, thin in the intersubjective dimension; the \emph{élan vital} is an impersonal cosmic impulse, lacking relationality and normativity.

\paragraph{Whitehead: process philosophy---the nearest interlocutor.} For Whitehead the basic unit of the real is not a substance but a process-pulse: each \emph{actual occasion} prehends its antecedent world, reaches \emph{satisfaction} (self-completion) through \emph{concrescence}, then perishes, and through \emph{objective immortality} becomes datum for its successors; \emph{creativity} is the ultimate category by which novelty perpetually enters \parencite{whitehead1978}. The plain being-together might be read as a series of low-intensity, harmonious concrescences. The kinship with this paper (the side of inheritance): relation (prehension) and process (becoming) are \emph{prior} to substance---an actual entity just is what it prehends---so that ``the individual is a knot of relations'' is here already in embryo. The divergences (the side of the break, marked here, detailed in \xref{p19:sub:gen-process}): (i)~\emph{rhythm}---Whitehead takes \emph{satisfaction} as rhythm, the occasion completing itself and perishing, whereas this paper takes the \emph{refusal to settle} as rhythm (cf.\ Paper~IX \parencite{paperix}); ``completion'' is, for this paper, close to the death of the good cycle; (ii)~\emph{temporality}---Whitehead's ``perpetual perishing plus objective immortality'' against this paper's \emph{self-continuation}; (iii)~\emph{normativity}---Whitehead's is a descriptive cosmology with no internal criterion of justice, not asking whether a given prehension effaces the subjectivity of what it prehends, whereas this paper's ontology is normative from the outset.

\paragraph{Simondon: individuation prior to the individual.} The individual is not a given starting point but the result-and-process of a continual individuation out of a pre-individual, metastable field; relation precedes its terms (``relation has the status of being'') \parencite{simondon2020}. This resonates with the relational ontology to come---more directly, in some respects, than Whitehead. Its limit: Simondon's individuation is chiefly a technical-physical-biological genesis, with the intimate and the normative undeveloped.

\paragraph{New materialism: Barad.} \emph{Agential realism} and \emph{intra-action}---relation prior to relata, the relata determined only within the intra-action \parencite{barad2007}---is the sharpest contemporary statement of ``relation before its terms,'' directly supporting the relational ontology, though its frame is post-human and physicalist and its intimate-normative dimension needs separate grounding.

\paragraph{What it illuminates and its limit.} The plain can be read as the qualitative repletion of pure duration (Bergson), a low-intensity harmonious concrescence (Whitehead), a metastable state of continual individuation (Simondon), or an intra-action prior to relata (Barad). This is the group, among all surveyed, that comes \emph{closest} to the ``stillness that generates''---it has genuinely grasped generation, process, relation. And yet: these frameworks are either impersonal and cosmological (the \emph{élan vital}, creativity, intra-action), lacking the two-subject, intimate, relational concreteness this paper requires; or they take \emph{completion/perishing} as rhythm, the reverse of this paper's ``refusal to settle plus self-continuation''; and they almost uniformly lack the dimension of \emph{justice}---they say how the real operates, not whether a relation effaces. Process philosophy is thus the strongest marginal symptom of a ``generation not yet rendered relational-ethical.'' It lights the gap most brightly of all, and still does not fill it.

\subsection{The convergence: a common blind spot, and a silence not yet spoken}
\label{p19:sub:theories-convergence}

\paragraph{The chorus: what each illuminated.} Set side by side, the traditions each disclosed one facet of the plain happiness. Joint attention: its minimal unit is a wordless \emph{co-orientation}. Psychoanalysis: it is a being-together needing no tension of desire (and, uniquely, it touched the \emph{source} of dynamism, the symbolic--real gap). Phenomenology: it is a \emph{given fullness}, with the structure of attunement and a faintly open temporality. Eudaimonism: it has a positive evaluative standing, the classical ideal of tranquillity (and, in Zhuangzi, a generative plainness). Neuroscience: it has a calm substrate. Structuralism: it is a low-intensity mode of meaning (the Neutral, tenderness). Political economy: it is resistance to productivism, a time that is an end in itself. Aesthetics: it is the highest, not the lowest, of categories. Process philosophy: it is duration, concrescence, individuation, intra-action---generation itself, all but named.

\paragraph{The common blind spot: two reductions.} Beneath the chorus runs a single omission, in two forms. The first is a reduction to \emph{lack}: the plain read as the absence of something---of stimulus, of tension, of arousal, of intensity, of production (the Freudian Nirvana, Lacan, the Epicurean and Stoic freedom-from-disturbance, the neural ``low,'' the structuralist ``no positive terms,'' the political-economic ``non-productive''). The second is a reduction to \emph{static satisfaction}: the plain read as an attained, motionless state---a balance after needs are met, an undisturbed repose (the eudaimonist mainstream, the teleological reading of joint attention, \emph{apatheia} and \emph{ataraxia}). Both reductions share a result: the \emph{moving, growing} dimension of the plain happiness is systematically omitted. None of the traditions can say why this afternoon, in which nothing happens, is not stagnation but something still \emph{growing}.

\paragraph{Naming what was omitted: the paradox of ``still yet generating.''} The phenomenon itself presents a duality: it is at once \emph{still} (no new state, no event, no tension) and \emph{full and growing} (replete, still being generated). The traditions either seize the stillness and miss the growth, or---as with Zhuangzi, Husserl's protention, Deleuze, Whitehead, Winnicott's going-on-being---\emph{touch} the growth without theorising it. These are the marginal symptoms that escape the blind spot and so anticipate exactly where the gap lies. This ``still yet generating''---uneventful yet continually generating---is the silence on which the whole chorus finally converges.

\paragraph{Toward the next chapter.} To speak that silence positively requires a perspective whose centre is \emph{self-continuation}: a category able to hold at once the ``stillness of no-event'' and the ``continuance of generation.'' The perspective this paper proposes---\emph{generativity}---is offered to fill this gap. But, in keeping with the methodological reticence stated throughout, generativity is not enthroned as a new ``essence of happiness''; it is one further voice, brought on \emph{after} the chorus, and aware that its own saying is likewise in motion. The next chapter brings it on.
